%% Cover letter -- Environmental Science & Technology (ACS ES&T), Research Article
%% Standalone document. In Overleaf: Menu -> Main document -> cover_letter.tex to compile it.
%% Items marked [TO ADD] need input from the authors before submission.
\documentclass[11pt,a4paper]{article}
\usepackage[margin=2.5cm]{geometry}
\usepackage[T1]{fontenc}\usepackage[utf8]{inputenc}\usepackage{lmodern}
\usepackage{xcolor,hyperref,enumitem}
\setlength{\parindent}{0pt}\setlength{\parskip}{0.7em}
\newcommand{\todo}[1]{\textcolor{red}{[TO ADD: #1]}}

\begin{document}

\textbf{To the ES\&T Editorial Team,}\\
Environmental Science \& Technology (ES\&T)

\bigskip

Dear ES\&T Editorial Team,

We are pleased to submit our manuscript, \textbf{``Beaver engineering creates intense, context-dependent dissolved organic carbon control points in streams''}, as a Research Article for consideration in \emph{Environmental Science \& Technology}.

\textbf{Environmental relevance.} Beavers are recolonizing Europe and North America, and their dams are increasingly used in river restoration. Dissolved organic carbon (DOC) is a key energy source for aquatic food webs, can deplete oxygen at high concentrations, and forms disinfection byproducts during drinking-water treatment. Inputs of organic matter are a leading cause of falling oxygen levels in receiving waters, and the form in which the matter arrives, dissolved or particulate, changes how strongly it affects water quality [1]. Whether beaver ponds raise or lower stream DOC is therefore a water-quality question with direct consequences for ecosystem management, yet earlier studies, mostly at single sites, reported conflicting results. Our study examined how beaver engineering changes DOC across 180 streams throughout Switzerland, sampled upstream and downstream of beaver ponds in winter and summer, and identified the conditions under which ponds act as DOC sources, sinks or neutral systems.

\textbf{Main findings, and why they matter.} We found that beaver ponds do not consistently raise or lower DOC. In some streams a pond adds DOC, in others it removes DOC, and in many it makes little difference, so the average effect looks close to zero while individual streams change a lot. Ponds raise DOC mainly by slowing the water, which allows more algae and plants to grow, and in summer they add an amount equal to about a fifth of the seasonal DOC change that streams already receive from upstream. This is new because it shows, with a national dataset, that earlier conflicting results are not noise but a real feature of how beavers shape streams, and it explains when each outcome is likely. It matters because it lets managers anticipate where and when beaver recolonization will change DOC and drinking-water quality.

\textbf{Fit with ES\&T.} The work links ecology, hydrology and biogeochemistry, and was carried out by a team of ecologists, geographers, hydrologists and biogeochemists together with national agencies (the Swiss Federal Office for the Environment and the national beaver agency). It addresses water quality and ecosystem management with a large, standardized dataset, in line with ES\&T's interest in integrative research with practical environmental relevance. Our study builds on earlier ES\&T work on beaver ponds and organic matter in water. A survey of 17 beaver impoundments in Quebec found that dissolved organic carbon and nutrients were two to four times higher in pond outlets than in inlets [2], which matches the DOC gains we see in some of our streams. A study of Swedish streams found that the effect of beaver ponds on methylmercury depended on the pond's colonization history, being strong only where ponds were newly flooded [3], in line with our finding that the effect of a pond depends on its setting rather than on the dam alone. Those studies focused on mercury or on a single region; ours addresses DOC itself across 180 streams and two seasons.

\textbf{Suggested editors.} \todo{preferred editors, if any; any non-preferred editor with the reason for the request}

We declare no conflicts of interest. We confirm that this manuscript is original, has not been published previously, and is not under review by any other journal or publishing outlet. All authors approve this submission.

Thank you for considering our manuscript. We appreciate your time and look forward to your review.

Sincerely,

\textbf{Dr.\ Leonardo Capitani} on behalf of all authors\\
Department of Aquatic Ecology, Eawag -- Swiss Federal Institute of Aquatic Science and Technology, Ueberlandstrasse 133, 8600 D\"ubendorf, Switzerland\\
Email: \href{mailto:leocapi07@gmail.com}{leocapi07@gmail.com}

\bigskip\hrule\bigskip

\textbf{Manuscript title:} Beaver engineering creates intense, context-dependent dissolved organic carbon control points in streams

\textbf{Authors, affiliations and contact information}

\begin{enumerate}[leftmargin=*,itemsep=0.3em]
\item \textbf{Leonardo Capitani}\textsuperscript{*} (corresponding author) -- WSL, Swiss Federal Institute for Forest, Snow and Landscape Research, Birmensdorf, Switzerland; Department of Aquatic Ecology, Eawag, Swiss Federal Institute of Aquatic Science and Technology, D\"ubendorf, Switzerland. \href{mailto:leocapi07@gmail.com}{leocapi07@gmail.com}
\item \textbf{Joshua R.\ Larsen} -- School of Geography, Earth and Environmental Sciences, University of Birmingham, UK. \href{mailto:j.larsen@bham.ac.uk}{j.larsen@bham.ac.uk}
\item \textbf{Annegret Larsen} -- Soil Geography and Landscape Group, Wageningen University \& Research, The Netherlands. \href{mailto:annegret.larsen@wur.nl}{annegret.larsen@wur.nl}
\item \textbf{Matthew Dennis} -- CIRCLE, Department of Geography, School of Environment and Life Sciences, The University of Manchester, UK. \href{mailto:matthew.dennis@manchester.ac.uk}{matthew.dennis@manchester.ac.uk}
\item \textbf{Lukas Hallberg} -- River Ecosystems Laboratory, Alpine and Polar Environmental Research Center, \'Ecole Polytechnique F\'ed\'erale de Lausanne, Sion, Switzerland; School of Geography, Earth and Environmental Sciences, University of Birmingham, UK. \href{mailto:lukas.hallberg1@gmail.com}{lukas.hallberg1@gmail.com}
\item \textbf{Valentin Moser} -- WSL, Swiss Federal Institute for Forest, Snow and Landscape Research, Birmensdorf, Switzerland; Department of Aquatic Ecology, Eawag, Swiss Federal Institute of Aquatic Science and Technology, D\"ubendorf, Switzerland. \href{mailto:valentin.moser@wsl.ch}{valentin.moser@wsl.ch}
\item \textbf{Christof Angst} -- Info fauna Nationale Biberfachstelle, Neuch\^atel, Switzerland. \href{mailto:christof.angst@infofauna.ch}{christof.angst@infofauna.ch}
\item \textbf{Silvan Minnig} -- umweltbildner.ch, Bern, Switzerland. \href{mailto:silvan.minnig@umweltbildner.ch}{silvan.minnig@umweltbildner.ch}
\item \textbf{Kaspar Berger} -- Institute of Geography, University of Bern, Hallerstrasse 12, CH-3012 Bern, Switzerland. \href{mailto:bergersumis@gmail.com}{bergersumis@gmail.com}
\item \textbf{Steffen Boch} -- WSL, Swiss Federal Institute for Forest, Snow and Landscape Research, Birmensdorf, Switzerland. \href{mailto:steffen.boch@wsl.ch}{steffen.boch@wsl.ch}
\item \textbf{Massimiliano Zappa} -- WSL, Swiss Federal Institute for Forest, Snow and Landscape Research, Birmensdorf, Switzerland. \href{mailto:massimiliano.zappa@wsl.ch}{massimiliano.zappa@wsl.ch}
\item \textbf{Bettina Schaefli} (Prof.) -- Institute of Geography, Hydrology, University of Bern, Hallerstrasse 12, CH-3012 Bern, Switzerland. \href{mailto:bettina.schaefli@unibe.ch}{bettina.schaefli@unibe.ch}
\item \textbf{Natalie Ceperley} (Dr.) -- Institute of Geography, University of Bern, Hallerstrasse 12, CH-3012 Bern, Switzerland. \href{mailto:natalie.ceperley@unibe.ch}{natalie.ceperley@unibe.ch}
\item \textbf{Francesco Pomati} -- Department of Aquatic Ecology, Eawag, Swiss Federal Institute of Aquatic Science and Technology, D\"ubendorf, Switzerland. \href{mailto:francesco.pomati@eawag.ch}{francesco.pomati@eawag.ch}
\item \textbf{Anita C.\ Risch} -- WSL, Swiss Federal Institute for Forest, Snow and Landscape Research, Birmensdorf, Switzerland. \href{mailto:anita.risch@wsl.ch}{anita.risch@wsl.ch}
\end{enumerate}

\bigskip\hrule\bigskip

\textbf{References}

\begin{enumerate}[leftmargin=*,itemsep=0.3em,label={[\arabic*]}]
\item McCabe, K.\,M.; Smith, E.\,M.; Lang, S.\,Q.; Osburn, C.\,L.; Benitez-Nelson, C.\,R. Particulate and dissolved organic matter in stormwater runoff influences oxygen demand in urbanized headwater catchments. \emph{Environ.\ Sci.\ Technol.} \textbf{2021}, 55, 952--961. \href{https://doi.org/10.1021/acs.est.0c04502}{doi:10.1021/acs.est.0c04502}
\item Roy, V.; Amyot, M.; Carignan, R. Beaver ponds increase methylmercury concentrations in Canadian Shield streams along vegetation and pond-age gradients. \emph{Environ.\ Sci.\ Technol.} \textbf{2009}, 43, 5605--5611.
\item Levanoni, O.; Bishop, K.; McKie, B.\,G.; Hartman, G.; Eklöf, K.; Ecke, F. Impact of beaver pond colonization history on methylmercury concentrations in surface water. \emph{Environ.\ Sci.\ Technol.} \textbf{2015}, 49, 12679. \href{https://doi.org/10.1021/acs.est.5b03146}{doi:10.1021/acs.est.5b03146}
\end{enumerate}

\end{document}
